\documentclass[11pt]{article}
\usepackage[T1]{fontenc}
\usepackage[utf8]{inputenc}
\usepackage{lmodern}
\usepackage[margin=1in]{geometry}
\usepackage{amsmath,amssymb,amsthm,mathtools}
\usepackage{enumitem}
\usepackage[hidelinks]{hyperref}

\newtheorem{theorem}{Theorem}[section]
\newtheorem{proposition}[theorem]{Proposition}
\newtheorem{definition}[theorem]{Definition}
\newtheorem{lemma}[theorem]{Lemma}
\newtheorem{corollary}[theorem]{Corollary}
\newtheorem{remark}[theorem]{Remark}
\newtheorem{example}[theorem]{Example}

\title{Sharp Support-Image Thresholds for Alternating Pencils}
\author{Luca Blanchi}
\date{June 2026}

\begin{document}
\maketitle

\begin{abstract}
We prove an exact arity threshold for structured support-image observers on a natural infinite family of alternating pencils and associated class-two exponent-$p$ groups. Let $q$ be odd. For every monic irreducible polynomial $f\in\mathbb F_q[t]$ of degree $m\ge3$, we construct a radical-free surjective alternating pencil $\beta_f:\Lambda^2 V_f\to W$ with $\dim_{\mathbb F_q}V_f=2m$ and $\dim_{\mathbb F_q}W=2$, whose Pfaffian is, up to scalar, the norm $N_{\mathbb F_{q^m}/\mathbb F_q}(a+b\theta)$, where $\theta$ is a root of $f$. Equivalently, $\beta_f$ is the regular irreducible symplectic pencil attached to $f$. For $r\ge1$, let $\operatorname{CSI}_{\le r}(\beta)$ denote the structured support-image profile through level $r$. We prove that $\operatorname{CSI}_{\le r}(\beta_f)$ is independent of $f$ for every $r<m$, and that the residual projective observer symmetry is the full group $PGL_2(q)$ at every such level. At level $m$, the profile detects $f$ up to the natural $PGL_2(q)$-action. Thus, on irreducible regular pencils of degree $m$, the exact support-image separation depth is $m$. For $q=p$ prime, the associated Baer groups are directly indecomposable special class-two exponent-$p$ groups of order $p^{2m+2}$, so the exact observer threshold is $(\log_p |G|-2)/2$.
\end{abstract}

\section{Introduction}

Alternating maps

\[
\beta:\Lambda^2V\to W
\]

control finite groups of nilpotency class two and exponent (p). If (p) is odd, such a map defines a group

\[
G_\beta=V\oplus W
\]

with multiplication

\[
(v,z)(u,w)
\left(
v+u,
z+w+\frac12\beta(v,u)
\right).
\]

The commutator is

\[
[(v,z),(u,w)]=(0,\beta(v,u)).
\]

Thus isomorphism questions for special class-two exponent-(p) groups are equivalent to pseudo-isometry questions for alternating maps.

This paper studies how much of an alternating map is visible to support-image observers.

For

\[
\mathbf x=(x_1,\ldots,x_r)\in V^r,
\]

define

\[
A_{\mathbf x}:V\to W^r,
\qquad
z\mapsto
\bigl(\beta(x_1,z),\ldots,\beta(x_r,z)\bigr).
\]

The level-(r) support-image profile records the image subspaces

\[
\operatorname{im}A_{\mathbf x}\le W^r
\]

as $\mathbf x$ varies. It is a structured invariant of $\beta$, taken up to the natural action of $GL(W)$.

The central question is:

\[
\boxed{
\text{How large must }r\text{ be before support-image data detects the map?}
}
\]

We answer this exactly for regular irreducible alternating pencils.

For each monic irreducible polynomial

\[
f\in\mathbb F_q[t]
\]

of degree (m), we construct a $2m$-dimensional alternating pencil

\[
\beta_f:\Lambda^2V_f\to\mathbb F_q^2.
\]

All nonzero scalar forms in the pencil are nondegenerate. Hence every first-order scalar-rank invariant sees the same object: a pencil in which every scalar member has full rank.

Nevertheless, the pencils are usually nonisomorphic. Their difference is encoded in the degree-(m) Pfaffian.

The main result says that support-image observers detect exactly this degree.

For every

\[
r<m,
\]

the complete structured support-image profile through level (r) is identical for all monic irreducible (f) of degree (m). Moreover, the residual projective symmetry visible to the observer is the full group

\[
PGL_2(q).
\]

At level (m), a single companion linearization of (f) becomes singular on $\beta_f$, and this detects (f) up to projective equivalence.

Thus the exact threshold is

\[
r=m.
\]

In group-theoretic terms, for $q=p$ prime, the associated groups have order

\[
p^{2m+2},
\]

so the exact threshold is

\[
\frac{\log_p |G|-2}{2}.
\]

This gives a linear observer lower bound with a matching upper bound, on an exponentially large family of directly indecomposable special groups.

\section{Structured support-image profiles}

Let (q) be odd, and let (V,W) be finite-dimensional $\mathbb F_q$-vector spaces.

An alternating map is a linear map

\[
\beta:\Lambda^2V\to W.
\]

Equivalently, it is an alternating bilinear map

\[
V\times V\to W.
\]

For $x\in V$, define the contraction

\[
A_x:V\to W,
\qquad
z\mapsto \beta(x,z).
\]

For an (r)-tuple

\[
\mathbf x=(x_1,\ldots,x_r)\in V^r,
\]

define

\[
A_{\mathbf x}:V\to W^r,
\qquad
z\mapsto
(A_{x_1}z,\ldots,A_{x_r}z).
\]

The level-(r) support-image profile is the function

\[
T\le W^r
\longmapsto
E_r^\beta(T),
\]

where

\[
E_r^\beta(T)
\#\{\mathbf x\in V^r:\operatorname{im}A_{\mathbf x}=T\}.
\]

The structured level-(r) support-image profile is this function on the full subspace lattice of $W^r$, considered together with the natural diagonal action of $GL(W)$.

It is often more convenient to use the equivalent rank transform.

Let

\[
U\le (W^r)^*.
\]

Choose a basis

\[
u_1,\ldots,u_s
\]

of (U). Each $u_j$ is a linear functional on $W^r$. Define

\[
\Psi_U^\beta:V^r\to (V^*)^s
\]

by

\[
\Psi_U^\beta(\mathbf x)
\bigl(
u_1\circ A_{\mathbf x},
\ldots,
u_s\circ A_{\mathbf x}
\bigr).
\]

Set

\[
\rho_r^\beta(U)
\operatorname{rank}\Psi_U^\beta.
\]

This is independent of the chosen basis of (U).

\begin{proposition}[CSI/rank-transform equivalence]
For fixed (r), the structured support-image profile

\[
T\mapsto E_r^\beta(T)
\]

and the rank-transform function

\[
U\mapsto \rho_r^\beta(U)
\]

determine one another.

\end{proposition}

\begin{proof}
For

\[
S\le W^r,
\]

define

\[
F_r^\beta(S)
\#\{\mathbf x\in V^r:\operatorname{im}A_{\mathbf x}\le S\}.
\]

The condition

\[
\operatorname{im}A_{\mathbf x}\le S
\]

is equivalent to

\[
u\circ A_{\mathbf x}=0
\qquad
\text{for all }u\in S^\perp.
\]

Thus

\[
F_r^\beta(S)
|\ker \Psi_{S^\perp}^\beta|
q^{r\dim V-\rho_r^\beta(S^\perp)}.
\tag{2.1}
\]

On the other hand,

\[
F_r^\beta(S)=\sum_{T\le S}E_r^\beta(T).
\tag{2.2}
\]

Mobius inversion on the lattice of subspaces of $W^r$ gives

\[
E_r^\beta(T)
\sum_{S\le T}
\mu(S,T)F_r^\beta(S).
\tag{2.3}
\]

Equations $2.1$--$2.3$ recover $E_r^\beta$ from $\rho_r^\beta$.

Conversely, $E_r^\beta$ determines $F_r^\beta$ by $2.2$, and then $2.1$ determines $\rho_r^\beta$.

$\square$

\end{proof}

\section{Residual projective observer symmetry}

Let
\[
PGL(W)=GL(W)/\mathbb F_q^\times.
\]

The group $PGL(W)$ acts on $W$, on $W^r$, and dually on $(W^r)^*$.

For

\[
g\in PGL(W)
\]

and

\[
U\le (W^r)^*,
\]

write

\[
gU
\]

for the image of (U) under the induced dual action.

\begin{definition}[Residual CSI symmetry]
The residual projective symmetry group of the structured CSI observer through level (r) is

\[
\Gamma_{\le r}(\beta)
\left\{
g\in PGL(W):
\rho_s^\beta(gU)=\rho_s^\beta(U)
\ \text{for all }1\le s\le r,\ U\le (W^s)^*
\right\}.
\]

Equivalently, by Proposition 2.1, $\Gamma_{\le r}(\beta)$ is the subgroup of $PGL(W)$ preserving all structured support-image profiles through level $r$.

If $(T,g)\in GL(V)\times GL(W)$ is a pseudo-isometry of $\beta$, i.e.

\[
g\beta(x,y)=\beta(Tx,Ty),
\]

then the projective class of (g) lies in $\Gamma_{\le r}(\beta$) for every (r). Thus $\Gamma_{\le r}(\beta$) may be larger than the true projective automorphism group: it measures the false symmetry not yet broken by the observer.

\end{definition}

\section{Grassmann image profiles}

We also use a parallel observer based on restrictions to subspaces.

For

\[
r\ge2
\]

and

\[
T\le W,
\]

define

\[
\mathcal E_r^\beta(T)
\#\left\{
L\le V:
\dim L=r,\
\beta(\Lambda^2L)=T
\right\}.
\]

Define also

\[
\mathcal F_r^\beta(U)
\#\left\{
L\le V:
\dim L=r,\
\beta(\Lambda^2L)\le U
\right\}.
\]

Then

\[
\mathcal F_r^\beta(U)
\sum_{T\le U}\mathcal E_r^\beta(T),
\]

and hence

\[
\mathcal E_r^\beta(T)
\sum_{U\le T}
\mu(U,T)\mathcal F_r^\beta(U),
\]

where $\mu$ is the Mobius function of the subspace lattice.

The condition

\[
\beta(\Lambda^2L)\le U
\]

means that (L) is totally isotropic for the quotient map

\[
\overline\beta_U:\Lambda^2V\to W/U.
\]

Thus the Grassmann image profile is equivalent to the collection of common isotropic-subspace counts for all quotients of the codomain.

This observer is related to, but not identical with, the tuple-based structured CSI profile. It is useful here because its two-plane collapse clarifies why level one is insufficient.

\section{The full two-plane profile collapses}

The case $r=2$ of the Grassmann image profile is special.

Let

\[
N_2=\binom{\dim V}{2}_q
\]

be the number of two-dimensional subspaces of (V).

Let

\[
z_\beta=\mathcal E_2^\beta(0)
\]

and, for

\[
P\in\mathbb P(W),
\]

let

\[
m_\beta(P)=\mathcal E_2^\beta(P).
\]

Thus $m_\beta(P$) counts two-planes whose image is the projective direction (P).

For

\[
0\ne a\in W^*,
\]

let

\[
J_\beta(a)
\#\{L\le V:\dim L=2,\ (a\circ\beta)|_L=0\}.
\]

If

\[
H_a=\mathbb P(\ker a)\subseteq\mathbb P(W),
\]

then

\[
J_\beta(a)=z_\beta+\sum_{P\in H_a}m_\beta(P).
\tag{5.1}
\]

Let

\[
d=\dim W,
\qquad
v=\frac{q^d-1}{q-1},
\qquad
k=\frac{q^{d-1}-1}{q-1}.
\]

Summing $5.1$ over all projective covectors $[a]\in\mathbb P(W^*$) gives

\[
\sum_{[a]}J_\beta(a)
kN_2+(v-k)z_\beta.
\]

Hence

\[
z_\beta
\frac{\sum_{[a]}J_\beta(a)-kN_2}{v-k}.
\tag{5.2}
\]

Once $z_\beta$ is known, $5.1$ gives the hyperplane sums of $m_\beta$. The point-hyperplane incidence matrix of projective space satisfies

\[
M^TM=q^{d-2}I+\lambda J,
\qquad
\lambda=\frac{q^{d-2}-1}{q-1}.
\]

Thus it is invertible over $\mathbb Q$, and the hyperplane sums recover $m_\beta$.

It remains to observe that $J_\beta(a$) is determined by the rank of the scalar alternating form

\[
a\circ\beta.
\]

Let (B) be an alternating form on an (n)-dimensional vector space, and suppose

\[
\operatorname{rank}B=2s.
\]

Then (B) has radical dimension $n-2s$. A two-plane is nonisotropic for (B) precisely when, for one, equivalently every, ordered basis ((x,y)), one has

\[
B(x,y)\ne0.
\]

The number of ordered pairs $(x,y)$ with $B(x,y)\ne0$ is

\[
(q^n-q^{n-2s})q^{n-1}(q-1).
\]

Dividing by

\[
|\operatorname{GL}_2(q)|
\]

gives the number of nonisotropic two-planes. Therefore

\[
J(B)
\binom n2_q
\frac{
(q^n-q^{n-2s})q^{n-1}(q-1)
}{
|\operatorname{GL}_2(q)|
}.
\tag{5.3}
\]

This is strictly decreasing in (s). Hence (J(B)) determines $\operatorname{rank}B$, and conversely.

\begin{theorem}[Two-plane collapse]
The full structured two-plane Grassmann image profile

\[
(z_\beta,m_\beta)
\]

and the marked scalar-rank function

\[
[a]\in\mathbb P(W^*)
\longmapsto
\operatorname{rank}(a\circ\beta)
\]

determine one another.

In particular, the full two-plane Grassmann image profile contains no more information than the level-one scalar-rank profile.

\end{theorem}

\section{Three-plane profiles are strictly stronger}

The collapse in Theorem 5.1 is special to two-planes.

We give two five-dimensional pencils over $\mathbb F_q$, (q) odd, with identical two-plane profiles but distinct common isotropic three-plane counts.

Let

\[
W=\mathbb F_q^2.
\]

\subsection{The Kronecker pencil}

Let

\[
V_K=X\oplus Y,
\qquad
\dim X=2,\quad \dim Y=3,
\]

with bases

\[
x_1,x_2
\]

for (X), and

\[
y_0,y_1,y_2
\]

for (Y). Define

\[
\beta_K:\Lambda^2V_K\to W
\]

by

\[
\beta_K(x_1,y_0)=\beta_K(x_2,y_1)=(1,0),
\]

\[
\beta_K(x_1,y_1)=\beta_K(x_2,y_2)=(0,1),
\]

and all remaining basis pairings zero.

For $(a,b)\ne(0,0)$, the scalar form

\[
a\beta_{K,0}+b\beta_{K,1}
\]

has cross-block matrix

\[
\begin{pmatrix}
a&b&0\
0&a&b
\end{pmatrix}.
\]

This matrix has rank $2$, so the scalar alternating form has rank $4$. Thus every nonzero scalar form in the pencil has rank $4$.

The subspace (Y) is common totally isotropic. It is the only common totally isotropic three-space.

Indeed, suppose $L\le V_K$ is common totally isotropic with $\dim L=3$. Let

\[
a=\dim \operatorname{pr}_X(L).
\]

If $a=0$, then $L=Y$.

If $a=1$, then $\dim(L\cap Y)\ge2$. But for a nonzero

\[
x=\alpha x_1+\beta x_2,
\]

the map

\[
Y\to W,
\qquad
y\mapsto \beta_K(x,y)
\]

has matrix

\[
\begin{pmatrix}
\alpha&\beta&0\
0&\alpha&\beta
\end{pmatrix},
\]

which has kernel of dimension $1$. Hence $L\cap Y$ cannot have dimension $2$.

If $a=2$, then $L\cap Y\ne0$. But the common right radical in (Y) of (X) is zero. This is again impossible.

Therefore

\[
\mathcal E_3^{\beta_K}(0)=1.
\tag{6.1}
\]

\subsection{The extension-field pencil}

Let

\[
E=\mathbb F_{q^2}
\]

and let

\[
V_E=\mathbb F_q r\oplus E^2.
\]

View

\[
W=E
\]

as a two-dimensional $\mathbb F_q$-space. Define

\[
\beta_E((c,u),(c',v))=\det_E(u,v).
\]

Every nonzero scalar projection

\[
E\to\mathbb F_q
\]

composed with the determinant form is nondegenerate on $E^2$ and has radical exactly $\mathbb F_qr$. Hence every nonzero scalar form in the pencil again has rank $4$.

A common totally isotropic three-space must contain (r). Its projection to $E^2$ must be an (E)-line. Conversely, every (E)-line gives such a three-space. Therefore

\[
\mathcal E_3^{\beta_E}(0)
|\mathbb P^1(E)|
q^2+1.
\tag{6.2}
\]

\begin{corollary}
The pencils $\beta_K$ and $\beta_E$ have identical complete two-plane image profiles but distinct three-plane image profiles:

\[
\mathcal E_3^{\beta_K}(0)=1,
\qquad
\mathcal E_3^{\beta_E}(0)=q^2+1.
\]

Thus three-plane image data strictly refines two-plane image data.

\end{corollary}

\section{Irreducible symplectic pencils}

We now construct the main family.

Let

\[
f(t)\in\mathbb F_q[t]
\]

be monic irreducible of degree

\[
m\ge3.
\]

Let

\[
E=\mathbb F_q[t]/(f),
\qquad
\theta=t\bmod f.
\]

Let

\[
V=E^2
\]

viewed as a $2m$-dimensional $\mathbb F_q$-space.

Let

\[
W=\mathbb F_q^2.
\]

Define

\[
\delta:E^2\times E^2\to E
\]

by

\[
\delta(x,y)=x_1y_2-x_2y_1.
\]

Define

\[
\beta_f:V\times V\to W
\]

by

\[
\beta_f(x,y)
\left(
\operatorname{Tr}{E/\mathbb F_q}\delta(x,y),
\operatorname{Tr}{E/\mathbb F_q}(\theta\delta(x,y))
\right).
\tag{7.1}
\]

Let $B_0,B_1$ be the two coordinate forms.

For $(a,b)\ne(0,0)$,

\[
aB_0+bB_1
\operatorname{Tr}_{E/\mathbb F_q}
\bigl((a+b\theta)\delta(-,-)\bigr).
\]

Since $a+b\theta\ne0$, the trace pairing is nondegenerate and $\delta$ is symplectic over (E). Hence every nonzero scalar form in the pencil is nondegenerate over $\mathbb F_q$.

Consequently, $\beta_f$ is radical-free and surjective.

The Pfaffian of the pencil is

\[
\operatorname{Pf}(aB_0+bB_1)
c\cdot N_{E/\mathbb F_q}(a+b\theta)
\tag{7.2}
\]

for some nonzero scalar $c\in\mathbb F_q^\times$.

Thus the binary Pfaffian form is the homogenization of (f), up to scalar.

We use the following standard canonical-form input for alternating pencils.

Canonical-form input 7.1

Over a field of characteristic not $2$, pairs of skew-symmetric matrices under congruence admit a canonical decomposition into regular and singular summands. In the regular skew-symmetric part, irreducible summands are determined by their elementary divisors. For a regular irreducible alternating pencil, the pseudo-isometry class is determined, up to projective change of pencil coordinates, by the projective equivalence class of the irreducible elementary divisor.

This is the standard Kronecker--Sergeichuk classification for pairs of skew-symmetric forms under congruence. We use it only in this regular irreducible case.

\section{Coordinate form of CSI for pencils}

Identify

\[
W^*
\]

with the two-dimensional space of linear polynomials

\[
\mathcal L={a+bt:a,b\in\mathbb F_q}.
\]

Let

\[
U\le (W^*)^r.
\]

Choose a basis of (U). We represent it by an $s\times r$ matrix

\[
C_U(t)
\]

whose entries are linear polynomials in (t).

Let (A) denote multiplication by $\theta$ on $V=E^2$. Using the nondegenerate form $B_0$ to identify (V) with $V^*$, the map defining

\[
\rho_r^{\beta_f}(U)
\]

is represented by

\[
C_U(A):V^r\to V^s.
\]

Therefore

\[
\rho_r^{\beta_f}(U)
\operatorname{rank}_{\mathbb F_q} C_U(A).
\tag{8.1}
\]

More explicitly, if

\[
C_U(t)=(c_{ij}(t)),
\]

then $C_U(A$) is the block map whose ((i,j))-block is multiplication by

\[
c_{ij}(\theta)
\]

on $E^2$.

This formula is the bridge between support-image profiles and polynomial matrix rank.

\section{Complete blindness below degree (m)}

\begin{lemma}
Let (f) be irreducible of degree (m). If

\[
r<m,
\]

then

\[
\rho_r^{\beta_f}(U)
2m\cdot
\operatorname{rank}_{\mathbb F_q(t)}C_U(t).
\tag{9.1}
\]

In particular, the value is independent of (f).

\end{lemma}

\begin{proof}
Let

\[
k=\operatorname{rank}_{\mathbb F_q(t)}C_U(t).
\]

Take the Smith normal form of $C_U(t)$ over $\mathbb F_q[t]$:

\[
P(t)C_U(t)Q(t)
\operatorname{diag}
(s_1(t),\ldots,s_k(t),0,\ldots,0),
\]

where (P(t),Q(t)) are unimodular.

The product

\[
s_1(t)\cdots s_k(t)
\]

divides every nonzero $k\times k$ minor. Since the entries of $C_U(t)$ have degree at most $1$, every such minor has degree at most $k\le r$. Hence

\[
\deg s_i(t)<m
\]

for each nonzero invariant factor $s_i$.

Because (f) is irreducible of degree (m), every nonzero $s_i(t$) is coprime to (f(t)). Therefore

\[
s_i(A)
\]

is invertible on $V=E^2$.

The evaluated matrices

\[
P(A),Q(A)
\]

are invertible because (P,Q) are unimodular.

Thus $C_U(A$) has exactly (k) invertible (V)-blocks and rank

\[
k\dim_{\mathbb F_q}V=2mk.
\]

$\square$

We now verify the projective invariance used later.

\end{proof}

\begin{lemma}[Projective invariance below threshold]
Let

\[
g\in PGL_2(q).
\]

For every $U\le (W^*)^r$,

\[
\operatorname{rank}_{\mathbb F_q(t)} C_{gU}(t)
=
\operatorname{rank}_{\mathbb F_q(t)} C_U(t).
\]

Consequently, if $r<m$, then

\[
\rho_r^{\beta_f}(gU)=\rho_r^{\beta_f}(U).
\]

\end{lemma}

\begin{proof}
Choose a representative

\[
g(t)=\frac{\alpha t+\beta}{\gamma t+\delta}.
\]

The induced action on linear pencil coordinates sends a linear form in (t) to a nonzero scalar multiple of

\[
(\gamma t+\delta)\ell(g(t)).
\]

Thus, after clearing a common denominator, $C_{gU}(t$) is obtained from $C_U(t)$ by:

\end{proof}

\begin{itemize}
\item applying the field automorphism of $\mathbb F_q(t)$ induced by the fractional-linear substitution $t\mapsto g(t$);
\end{itemize}

\begin{itemize}
\item multiplying rows or columns by nonzero rational functions;
\end{itemize}

\begin{itemize}
\item changing bases in (U) and in $(W^*)^r$, i.e. applying invertible row and column operations over $\mathbb F_q$.
\end{itemize}

All these operations preserve rank over the rational function field $\mathbb F_q(t)$. Therefore

\[
\operatorname{rank}
{\mathbb F_q(t)} C
{gU}(t)
\operatorname{rank}_{\mathbb F_q(t)} C_U(t).
\]

The final statement follows from Lemma 9.1.

$\square$

\begin{theorem}[Maximal low-level symmetry]
For every monic irreducible (f) of degree (m),

\[
\Gamma_{\le r}(\beta_f)=PGL_2(q)
\qquad
\text{for every }r<m.
\tag{9.2}
\]

Moreover, if (f,h) are any two monic irreducible polynomials of degree (m), then

\[
\operatorname{CSI}_{\le m-1}(\beta_f)
\operatorname{CSI}_{\le m-1}(\beta_h).
\tag{9.3}
\]

Thus all support-image observers of arity below (m) are completely blind to the polynomial (f).

\end{theorem}

\begin{proof}
For $r<m$, Lemma 9.1 gives

\[
\rho_r^{\beta_f}(U)
2m\cdot
\operatorname{rank}_{\mathbb F_q(t)}C_U(t),
\]

which is independent of (f). This proves $9.3$.

Lemma 9.2 shows that every projective change of pencil coordinates preserves the rank transform through level $r<m$. Hence

\[
PGL_2(q)\le \Gamma_{\le r}(\beta_f).
\]

The reverse inclusion is automatic from the definition, since $\Gamma_{\le r}(\beta_f)$ is a subgroup of $PGL(W)=PGL_2(q)$. Therefore

\[
\Gamma_{\le r}(\beta_f)=PGL_2(q).
\]

$\square$

\end{proof}

\section{Level (m) detects the polynomial}

Let $C_f$ be the companion matrix of (f), and define the $m\times m$ linear matrix polynomial

\[
L_f(t)=tI_m-C_f.
\]

Then

\[
\det L_f(t)=f(t).
\]

Let

\[
U_f\le (W^*)^m
\]

be the row space of $L_f(t$).

For the pencil $\beta_f$, evaluation at (A) gives

\[
L_f(A),
\]

which is singular because

\[
f(A)=0.
\]

Thus

\[
\rho_m^{\beta_f}(U_f)<2m^2.
\tag{10.1}
\]

Now let (h) be another monic irreducible polynomial of degree (m). Let $A_h$ be multiplication by a root of (h) on $E_h^2$.

For

\[
g\in PGL_2(q),
\]

the transformed subspace

\[
gU_f
\]

is represented by a linear matrix polynomial whose determinant is, up to a nonzero scalar, the projective transform

\[
g\cdot f.
\]

If (h) is not projectively equivalent to (f), then (h) and $g\cdot f$ are distinct irreducibles of degree (m), hence coprime. Therefore

\[
(g\cdot f)(A_h)
\]

is invertible, and the evaluated linear matrix is invertible. Hence

\[
\rho_m^{\beta_h}(gU_f)=2m^2.
\tag{10.2}
\]

Equations $10.1$ and $10.2$ give a level-(m) witness separating $\beta_f$ from every projectively inequivalent $\beta_h$.

\begin{theorem}[Exact level-(m) classification]
Let $f,h\in\mathbb F_q[t]$ be monic irreducible polynomials of degree $m\ge3$. Then

\[
\operatorname{CSI}_{\le m}(\beta_f)
\operatorname{CSI}_{\le m}(\beta_h)
\]

up to the natural $PGL_2(q)$-action if and only if

\[
h\in PGL_2(q)\cdot f.
\]

Equivalently, $\beta_f$ and $\beta_h$ are pseudo-isometric as alternating pencils.

\end{theorem}

\begin{proof}
If $h$ is not projectively equivalent to $f$, the companion linearization above gives an explicit $U_f\le (W^*)^m$ such that the rank value for $\beta_f$ is singular, while every projective translate has full rank for $\beta_h$. Thus the level-$m$ CSI profiles differ.

Conversely, if $h$ is projectively equivalent to $f$, then the corresponding Pfaffian binary forms are projectively equivalent. By the canonical-form input for regular irreducible skew-symmetric pencils, the pencils $\beta_f$ and $\beta_h$ are pseudo-isometric. Hence all CSI profiles agree.

$\square$

Combining Theorems 9.3 and 10.1, the exact CSI separation depth of the irreducible degree-(m) pencil family is

\[
\boxed{m.}
\]

\end{proof}

\section{True and false projective symmetry}

Let

\[
\operatorname{Aut}_W(\beta_f)
\le PGL_2(q)
\]

be the projective image of the pseudo-isometry group of $\beta_f$.

The Pfaffian form implies

\[
\operatorname{Aut}_W(\beta_f)
\le
\operatorname{Stab}_{PGL_2(q)}(f).
\]

The converse follows from the regular symplectic pencil classification: a projective transformation stabilizing the elementary divisor is realized by a pseudo-isometry of the pencil. Hence

\[
\operatorname{Aut}_W(\beta_f)
\operatorname{Stab}_{PGL_2(q)}(f).
\tag{11.1}
\]

From Theorem 9.3,

\[
\Gamma_{\le r}(\beta_f)=PGL_2(q)
\qquad
(r<m).
\]

From Theorem 10.1,

\[
\Gamma_{\le m}(\beta_f)=\operatorname{Stab}_{PGL_2(q)}(f).
\tag{11.2}
\]

Therefore support-image observers exhibit maximal false symmetry below the threshold:

\[
\boxed{
\Gamma_{\le r}(\beta_f)=PGL_2(q)
\quad (r<m),
}
\]

but at the threshold the false symmetry collapses exactly to the true projective automorphism group:

\[
\boxed{
\Gamma_{\le m}(\beta_f)=\operatorname{Aut}_W(\beta_f).
}
\]

\section{Direct indecomposability of the associated groups}

Assume in this section that

\[
q=p
\]

is an odd prime. Then $G_{\beta_f}$ is a finite (p)-group of class two and exponent (p).

Since $\beta_f$ is radical-free and surjective,

\[
Z(G_{\beta_f})=G_{\beta_f}'=W.
\]

Thus $G_{\beta_f}$ is special.

\begin{proposition}
The group $G_{\beta_f}$ is directly indecomposable.

\end{proposition}

\begin{proof}
It is enough to show that the alternating pencil admits no nontrivial orthogonal direct decomposition.

Suppose

\[
V=V_1\perp V_2
\]

is a nontrivial decomposition orthogonal for both $B_0$ and $B_1$.

Since $B_0$ is nondegenerate, each $V_i$ is $B_0$-nondegenerate.

Let (A) be multiplication by $\theta$ on $V=E^2$. Then

\[
B_1(x,y)=B_0(Ax,y).
\]

Orthogonality for $B_1$ implies that each $V_i$ is (A)-invariant. Since (A) has irreducible minimal polynomial (f), (A)-invariant $\mathbb F_q$-subspaces of $E^2$ are precisely (E)-subspaces.

A proper nonzero (E)-subspace of $E^2$ is an (E)-line. Every (E)-line is totally isotropic for

\[
\delta(x,y)=x_1y_2-x_2y_1,
\]

and hence for $B_0$. This contradicts the $B_0$-nondegeneracy of $V_i$.

Therefore no nontrivial orthogonal direct decomposition exists, and $G_{\beta_f}$ is directly indecomposable.

$\square$

The order of the group is

\[
|G_{\beta_f}|=p^{2m+2}.
\]

Thus the exact threshold can be written as

\[
\boxed{
m=\frac{\log_p |G_{\beta_f}|-2}{2}.
}
\]

\end{proof}

\section{Exponentially large low-level fibres}

Let

\[
I_q(m)
\]

denote the number of monic irreducible polynomials of degree (m) over $\mathbb F_q$. Then

\[
I_q(m)
\frac1m
\sum_{e\mid m}\mu(e)q^{m/e}
\frac{q^m}{m}+O_q(q^{m/2}).
\]

The group $PGL_2(q)$ has order

\[
|PGL_2(q)|=q(q^2-1).
\]

Therefore the number of $PGL_2(q)$-orbits of irreducible degree-(m) polynomials is at least

\[
\frac{I_q(m)}{|PGL_2(q)|}
\frac{q^m}{m,q(q^2-1)}
+
O_q(q^{m/2}).
\tag{13.1}
\]

Choosing one representative from each orbit gives a family

\[
\mathcal F_m
\]

of pairwise non-pseudo-isometric irreducible pencils.

For $q=p$ prime, these yield pairwise nonisomorphic special directly indecomposable class-two exponent-(p) groups.

Every member of $\mathcal F_m$ has the same CSI profile through level

\[
m-1.
\]

At level (m), the profile separates the family.

Thus a single $\operatorname{CSI}_{\le m-1}$-fibre contains

\[
q^{m-o(m)}
\]

pairwise nonisomorphic objects. In group order,

\[
q^m=q^{(\log_q |G|-2)/2}=|G|^{1/2}q^{-1}.
\]

So the low-level fibre has size

\[
|G|^{1/2-o(1)}.
\]

\section{Generic triviality at the threshold}

The threshold often destroys all projective symmetry.

\begin{proposition}
Suppose

\[
\gcd(m,|PGL_2(q)|)=1.
\]

Then every monic irreducible polynomial of degree (m) has trivial $PGL_2(q)$-stabilizer.

\end{proposition}

\begin{proof}
Let $\alpha$ be a root of (f). The roots of (f) are the Frobenius orbit

\[
\alpha,\alpha^q,\ldots,\alpha^{q^{m-1}}.
\]

If

\[
1\ne g\in PGL_2(q)
\]

stabilizes (f), then (g) permutes this Frobenius orbit. Hence for some (j),

\[
g\alpha=\alpha^{q^j}.
\]

The induced permutation on the orbit has order dividing (m). On the other hand, the order of (g) divides

\[
|PGL_2(q)|.
\]

Since

\[
\gcd(m,|PGL_2(q)|)=1,
\]

the induced permutation must be trivial. Thus (g) fixes $\alpha$. But a nonidentity projective linear transformation over $\mathbb F_q$ has at most two fixed points over the algebraic closure, and cannot fix a degree-(m) element with $m\ge3$. Therefore $g=1$.

$\square$

There are infinitely many such (m). For these degrees,

\[
\Gamma_{\le m-1}(\beta_f)=PGL_2(q),
\qquad
\Gamma_{\le m}(\beta_f)=1.
\]

A probabilistic version also holds.

\end{proof}

\begin{proposition}
Let (f) be uniformly random among monic irreducible polynomials of degree (m). Then

\[
\Pr\left(
\operatorname{Stab}_{PGL_2(q)}(f)\ne1
\right)
q^{-\Omega_q(m)}.
\]

More explicitly,

\[
\Pr\left(
\operatorname{Stab}_{PGL_2(q)}(f)\ne1
\right)
O_q\left(q^{-m/(q+1)}\right).
\]

\end{proposition}

\begin{proof}
It is equivalent to sample a degree-(m) element $\alpha\in\overline{\mathbb F}_q$, up to Frobenius orbit.

Fix

\[
1\ne g\in PGL_2(q)
\]

of order (h). Every element of $PGL_2(q)$ has order at most $q+1$, after including the unipotent case since $p\le q<q+1$.

If (g) stabilizes the Frobenius orbit of a degree-(m) element $\alpha$, then

\[
g\alpha=\alpha^{q^j}
\]

for some $0\le j<m$.

If $j=0$, then (g) fixes $\alpha$. A nonidentity projective linear transformation has at most two fixed points over the algebraic closure, so this cannot occur for a degree-(m) element with $m\ge3$.

Thus $1\le j<m$. Applying (g) repeatedly gives

\[
\alpha=g^h\alpha=\alpha^{q^{jh}},
\]

so

\[
m\mid jh.
\]

Hence (j) is a multiple of $m/\gcd(m,h)$, and there are at most $h-1\le q$ possible values of (j).

For fixed (j), the equation

\[
g(x)=x^{q^j}
\]

has at most $q^j+1$ solutions in the algebraic closure. The largest possible (j) is at most

\[
m\left(1-\frac1h\right)
\le
m\left(1-\frac1{q+1}\right).
\]

Thus the number of degree-(m) elements whose Frobenius orbit is stabilized by this fixed (g) is

\[
O_q\left(q^{m(1-1/(q+1))}\right).
\]

The number of degree-(m) elements over $\mathbb F_q$ is

\[
q^m+O_q(q^{m/2}).
\]

Therefore, for fixed $g\ne1$, the probability that (g) stabilizes the orbit is

\[
O_q\left(q^{-m/(q+1)}\right).
\]

Taking a union bound over the finite group $PGL_2(q)$ gives the claim.

$\square$

Thus a random member of the irreducible pencil family has trivial true projective automorphism group with exponentially high probability, while its support-image observer symmetry is still the full $PGL_2(q)$ below level (m).

\end{proof}

\section{Main theorem}

We collect the results.

\begin{theorem}[Sharp support-image threshold for irreducible alternating pencils]
Let (q) be odd and $m\ge3$. For every monic irreducible polynomial $f\in\mathbb F_q[t]$ of degree (m), let

\[
\beta_f:\Lambda^2\mathbb F_q^{2m}\to \mathbb F_q^2
\]

be the alternating pencil defined in $7.1$.

Then:

\end{theorem}

\begin{itemize}
\item For every $r<m$, the structured support-image profile $\operatorname{CSI}_{\le r}(\beta_f)$ is independent of (f).
\end{itemize}

\begin{itemize}
\item For every $r<m$,
\end{itemize}

\[
\Gamma_{\le r}(\beta_f)=PGL_2(q).
\]

\begin{itemize}
\item At level (m), the profile detects (f) up to the natural $PGL_2(q)$-action.
\end{itemize}

\begin{itemize}
\item Hence the exact CSI separation depth of the irreducible degree-(m) pencil family is (m).
\end{itemize}

\begin{itemize}
\item A single $\operatorname{CSI}_{\le m-1}$-fibre contains at least
\end{itemize}

\[
\frac{q^m}{m,q(q^2-1)}
+
O_q(q^{m/2})
\]

distinct pseudo-isometry classes.

\begin{itemize}
\item For a proportion
\end{itemize}

\[
1-O_q(q^{-m/(q+1)})
\]

of irreducible (f), the level-(m) residual observer symmetry is trivial:

\[
\Gamma_{\le m}(\beta_f)=1.
\]

If $q=p$ is prime, the corresponding Baer groups $G_{\beta_f}$ are directly indecomposable special class-two exponent-(p) groups of order

\[
p^{2m+2},
\]

and the exact observer threshold is

\[
m=\frac{\log_p |G_{\beta_f}|-2}{2}.
\]

\begin{proof}
Parts $1$ and $2$ are Theorem 9.3. Parts $3$ and $4$ are Theorem 10.1. Part $5$ is the orbit count in Section 13. Part $6$ is Proposition 14.2 together with $11.2$. The group-theoretic statement follows from Proposition 12.1 and the order calculation.

$\square$

\end{proof}

\section{Consequences for star, flag, and rank-support profiles}

The support-image profile dominates the rank-support and flag profiles.

For

\[
\mathbf x\in V^r,
\]

the rank of

\[
A_{\mathbf x}:V\to W^r
\]

is the dimension of the support image

\[
\operatorname{im}A_{\mathbf x}.
\]

Thus CSI determines the rank distribution

\[
N_r^\beta(a)
\#\{\mathbf x\in V^r:\operatorname{rank}A_{\mathbf x}=a\}.
\]

It also determines the aggregated flag profile, because for

\[
S\subseteq{1,\ldots,r},
\]

the rank of the coordinate subtuple is the dimension of the projection of

\[
\operatorname{im}A_{\mathbf x}\le W^r
\]

to the coordinates in (S).

Therefore Theorem 15.1 implies:

\begin{corollary}
For the family

\[
{\beta_f:\deg f=m},
\]

all star profiles, rank-support profiles, and aggregated flag profiles of arity

\[
r<m
\]

are identical.

At arity (m), structured CSI separates the irreducible-pencil family exactly.

This shows that the obstruction is not merely a failure of crude commuting probabilities; it persists through the full structured support-image hierarchy below the threshold.

\end{corollary}

\section{Interpretation}

The theorem identifies a precise mechanism behind support-image depth.

For an irreducible degree-(m) pencil, all polynomial tests of matrix size below (m) see only generic rank over the rational function field

\[
\mathbb F_q(t).
\]

The irreducible polynomial (f) cannot divide any nonzero determinant of degree $<m$. Thus every low-level observer is blind to (f).

At level (m), a determinant of degree (m) can equal (f). The companion linearization

\[
tI_m-C_f
\]

is the minimal witness. Its determinant is precisely (f), and it becomes singular exactly on the corresponding pencil.

Thus the observer transition is governed by the first possible degree of a Pfaffian/eigenvalue witness:

\[
\boxed{
\text{support-image depth}
\text{degree of the irreducible elementary divisor}.
}
\]

Below threshold, the observer sees only rational-function-field rank. At threshold, it sees the finite-field point $\theta$.

\section{Further directions}

\subsection{Higher-dimensional regular spaces}

The present paper treats pencils, i.e. $\dim W=2$. The natural next case is a regular alternating space of derived rank $d\ge3$. One expects similar thresholds governed by the minimal degree of determinantal or Pfaffian equations cutting out the spectral data.

\subsection{Generic support-image rigidity}

The result suggests that generic random alternating spaces may have an observer threshold governed by the minimal degree of their projective degeneracy locus. Proving this would require anti-concentration for random Pfaffian systems.

\subsection{Deterministic Grassmannian projection asymmetry}

The full deterministic two-plane profile collapses to scalar ranks, but higher-plane image profiles may define asymmetric projective fibre counts for random projections of Grassmannians. This remains an attractive route to generic deterministic support-image canonization.

\subsection{Algorithms}

The level-(m) witness is explicit: a companion linearization. This suggests an algorithmic strategy for regular pencils: search for the first arity (r) at which a support-image rank defect appears. The first defect reveals the degree of the irreducible elementary divisor.


\section*{Declaration of generative AI and AI-assisted technologies in the writing process}
During the preparation of this work, ChatGPT, by OpenAI, was used to assist with mathematical drafting, formalization, review, and editing. This work is shared as a preliminary AI-assisted mathematical note. The mathematical content may have been only partially reviewed and may contain errors; it should not be treated as peer-reviewed or as a fully verified manuscript.

\begin{thebibliography}{99}
\bibitem{ref1} V. A. Bovdi, T. G. Gerasimova, M. A. Salim, and V. V. Sergeichuk. Reduction of a pair of skew-symmetric matrices to its canonical form under congruence. Linear Algebra and its Applications 498 $2016$, 403--433.
\bibitem{ref2} R. Gow and G. McGuire. Invariant rational functions, linear fractional transformations and irreducible polynomials over finite fields. Finite Fields and Their Applications 78 $2022$, 101977.
\bibitem{ref3} R. A. Horn and V. V. Sergeichuk. Canonical forms for complex matrix congruence and ()-congruence*. Linear Algebra and its Applications 416 $2006$, 1010--1032.
\bibitem{ref4} R. Lidl and H. Niederreiter. Finite Fields. Cambridge University Press, 1997.
\bibitem{ref5} L. Reis. Invariant theory of a special group action on irreducible polynomials over finite fields. Designs, Codes and Cryptography 87 $2019$, 1095--1117.
\bibitem{ref6} V. V. Sergeichuk. Canonical matrices of bilinear and sesquilinear forms. Linear Algebra and its Applications 343--344 $2002$, 43--60.
\bibitem{ref7} V. V. Sergeichuk. Canonical matrices of forms and pairs of forms over finite and (p)-adic fields. Linear Algebra and its Applications 438 $2013$, 2322--2339.
\bibitem{ref8} J. B. Wilson. Decomposing (p)-groups via Jordan algebras. Journal of Algebra 322 $2009$, 2642--2679.
\bibitem{ref9} P. A. Brooksbank and J. B. Wilson. Computing isometry groups of Hermitian maps. Transactions of the American Mathematical Society 364 $2012$, 1975--1996.
\bibitem{ref10} Y. Li and Y. Qiao. Linear algebraic analogues of the graph isomorphism problem and the Erdos--Renyi model. SIAM Journal on Computing 51 $2022$, 926--963.
\end{thebibliography}
\end{document}
